% Generated from authoritative CSVs by render_revision_artifacts.py.
\newcommand{\SystemDiff}{-3.00}
\newcommand{\SystemCILow}{-7.50}
\newcommand{\SystemCIHigh}{+1.35}
\newcommand{\PathFinalMAE}{0.1300}
\newcommand{\ExitFinalMAE}{0.0107}
\newcommand{\PathObservedPct}{22.17}
\newcommand{\ExitObservedPct}{97.33}
\newcommand{\DiagnosticResultsText}{Of 180 planned snapshots, 180 were eligible; 720 of 720 counterfactual branches completed. Unconditioned control selected a recommendation equivalent to reactive in 0.0\% of snapshots. Its fixed-action clearance ranking reversed in 28.9\% (95\% seed-bootstrap interval [20.0, 37.8]\%). With the common reactive continuation, the selected-minus-reactive action effect was -0.67~s ([-5.44, +4.44]~s). Conditioned control selected a recommendation equivalent to reactive in 4.4\% of snapshots. Its fixed-action clearance ranking reversed in 14.4\% (95\% seed-bootstrap interval [8.9, 20.0]\%). With the common reactive continuation, the selected-minus-reactive action effect was -3.56~s ([-7.56, +0.11]~s). Both action-effect intervals include zero. Reversal frequencies use all completed decisions as denominator; predicted or actual ties are not reversals and are separately tabulated. The diagnostics demonstrate that forecast ranking need not persist under realized capacity; they do not apportion the contributions of capacity mismatch, terminal approximation, finite scenarios and objective choice. The two reversal rates describe different on-policy state sets; their difference is not an isolated effect of conditioning.}
\newcommand{\RuntimeResultsText}{On AMD64 Family 26 Model 104 Stepping 0, AuthenticAMD, Python 3.14.6 and NumPy 2.5.3, the benchmark recorded 2700 timed decisions over 30 shared snapshots. Both predictive controllers had similar latency, with overall p95 at most 4.59~ms; reactive control was substantially cheaper. These measurements exclude state copying and I/O and provide no hard real-time or deployment guarantee.}
